\section{Algorithm and Hardware Co Design}
\label{sec:codesign}
The encoder has four pre-normalized blocks, width $d=128$, four attention heads,
and feed-forward width 256. It uses RMSNorm, piecewise-linear SiLU, learned
position embeddings, mean pooling, and classification heads at the final and
block-two boundaries. The current modes use eight-bit activations; four-bit
weights do not imply a DSP subword-packing speedup.

After block one, a learned score is reduced within chronological groups of four
tokens. Stable ranking retains the earlier group on an exact score tie. Each
retained group contributes four original tokens. Residual Token Summary (RTS)
adds one importance-weighted summary for every discarded group. With eight
input groups and $k$ retained groups, the later sequence length is
\begin{equation}
  N_{\mathrm{RTS}}=4k+(8-k).
\end{equation}
In real-valued form, a group's summary uses weights $|s_i|+1$; its implementation
uses the frozen fixed-point representation, rounding, and checked widths.
Compaction preserves chronological order.

\begin{table}
\caption{Current mode configuration. Token counts apply after the first block.}
\label{tab:modes}
\centering
\begin{tabular}{lccc}
\toprule
Mode & Weight bits & Later tokens & Executed blocks \\
\midrule
Quality & 8 & 32 & 4 \\
Balanced & 4 & 26 & 4 \\
Energy-Saver & 4 & 20 & 2 \\
\bottomrule
\end{tabular}
\end{table}

The frozen default uses independently trained modes. A separately published
shared HAR family uses one shadow parameter set per seed, with mode-specific
calibration and encoded constants. These families have different training
budgets and must not be compared as an isolated sharing experiment.
\writingtask{model lead}{Document training losses, checkpoint selection,
straight-through surrogates, calibration, and exact integer summary arithmetic.
Present dropping and RTS controls at both equal retention and equal processed
sequence length.}
